\documentclass[conference]{IEEEtran}
\IEEEoverridecommandlockouts
\usepackage{cite}
\usepackage{amsmath,amssymb,amsfonts}
\usepackage{algorithmic}
\usepackage{graphicx}
\usepackage{textcomp}
\usepackage{xcolor}
\usepackage{hyperref}
\usepackage{booktabs}
\usepackage{url}
\usepackage{stfloats}
\def\BibTeX{{\rm B\kern-.05em{\sc i\kern-.025em b}\kern-.08em
    T\kern-.1667\lower.7ex\hbox{E}\kern-.125emX}}

\hypersetup{colorlinks=true, linkcolor=blue, urlcolor=blue, citecolor=blue}

\begin{document}

\title{Linwarm: Telugu-First TokEval --- Small Dedicated Tokenizers Beat Massive\\ Multilingual Vocabularies for Phone-Scale LLMs}

\author{\IEEEauthorblockN{Gottipati Jamadagni}
\IEEEauthorblockA{\textit{Software Engineer, Unnanu} \\ \textit{Hyderabad, India}\\
\texttt{jama.gottipati7@gmail.com}}
\and
\IEEEauthorblockN{AI Collaborator (Muse Spark)}
\IEEEauthorblockA{\textit{Research Assistant}\\
\texttt{opencode / Muse Spark 1.2}}
}

\maketitle

\begin{abstract}
English-centric Byte-Pair Encoding (BPE) tokenizers impose a severe fertility penalty on morphologically rich, low-resource languages. On Telugu, \texttt{tiktoken cl100k} averages 13.57 tokens/word versus 1.14 on English (parity 11.9), inflating context length, latency, and cost for edge and cloud deployment. Massive multilingual vocabularies partially mitigate this: Sarvam 262K (2.48) and Saiteja 50K (1.64). Replicating the TokEval intrinsic evaluation suite \cite{meister2026tokeval} on Telugu, we train dedicated SentencePiece BPE and Unigram tokenizers on 80K Telugu Wikipedia lines (36 MB, \texttt{wikimedia/wikipedia 20231101.te}) and evaluate on FLORES-200 997-sentence parallel data. Dedicated 16K/32K vocabularies achieve \textbf{1.55 tokens/word} on Telugu --- \textbf{8.8$\times$ longer effective context than \texttt{tiktoken} and 37\% better than Sarvam 262K despite 8$\times$ smaller vocabulary}. Unigram better preserves morphological boundaries (stem+\textit{di} inflection split vs arbitrary split), consistent with recent Indic studies \cite{brahma2026indic}. A joint Telugu--English 16K BPE balances bilingual parity to 1.34 (TE 1.91, EN 1.42) with 64 MB embedding cost versus 2.1 GB for 262K, enabling 1.7B models to run on phones. We release 8 tokenizers, code, and reproducibility logs.
\end{abstract}

\begin{IEEEkeywords}
tokenization, Telugu, low-resource languages, TokEval, parity, edge LLM, SentencePiece, multilingual
\end{IEEEkeywords}

\section{Introduction}
The vision of 1--2B parameter language models that train, personalize, and infer directly on phones---learning patterns differently per user and language---is limited first by tokenization. A Spanish word typically requires 1--2 tokens with English-centric vocabularies, while an equivalent Telugu word requires 6--12, burning context window, memory, and inference budget \cite{meister2026tokeval, brahma2026indic, anylearn2026tokenization}. For a Hyderabad-based deployment serving Telugu and English on Azure and edge, this translates to 3--8$\times$ cost and latency disparity.

On-device training further amplifies the gap. A 1.7B model with a 100K vocabulary needs $\approx$400 MB embedding alone (vocab $\times$ hidden $\times$ 2 bytes), while a 262K vocabulary (Sarvam) needs 2.1 GB, exceeding typical 6 GB phone RAM after quantizing weights. A smaller, Telugu-aware vocabulary is not merely an efficiency tweak but an enabler for the ``different learning patterns'' paradigm: sparse, federated, and morphology-aware.

Recent work provides the analytical tools. TokEval \cite{meister2026tokeval} (COLM 2026) replaces expensive pretraining sweeps with intrinsic metrics---fertility, compression, UTF-8 boundary integrity, and digit place-value alignment---showing strong correlation ($\rho\approx0.80$) with bits-per-byte language modeling ability and task-specific accuracy. Brahma \textit{et al.} \cite{brahma2026indic} evaluate 17 Indic languages (11 scripts, Findings ACL 2026) and find (i) script-aware normalization helps, (ii) Unigram LM preserves morphological boundaries better than BPE, and (iii) cluster-based vocabulary construction beats joint training. Kanjirangat \textit{et al.} \cite{kanjirangat2026parity} show parity-aware BPE's fairness gains are script-sensitive and uneven. Yet no study replicates TokEval end-to-end on Telugu at Wikipedia scale with explicit phone-fit modeling and head-to-head comparison against deployed tokenizers (\texttt{tiktoken}, Sarvam, Saiteja).

This paper fills that gap. \textbf{Contributions:} (1) Large Telugu corpus (80K Wikipedia lines, 36 MB) plus FLORES-200 997-sentence parallel set for controlled evaluation; (2) 8 tokenizers: BPE vs Unigram at 16K/32K, Telugu-only vs joint Telugu--English vs simulated Dravidian cluster, with/without normalization; (3) Head-to-head TokEval comparison against \texttt{tiktoken cl100k} (100K), \texttt{sarvamai/sarvam-30b} 262K, and \texttt{Saiteja/telugu-bpe} 50K; (4) Phone embedding cost model and digit/UTF-8 failure analysis with next-step fixes; (5) Open release of tokenizers and reproducibility bundle.

\section{Related Work}

\subsection{Tokenizer Evaluation}
Standard metrics (fertility, compression) are insufficient \cite{meister2026tokeval}. TokEval adds UTF-8 integrity (whether a token splits a multi-byte character), digit place-value alignment (whether ``12345'' stays as one piece for arithmetic), and line-break handling, linking each to downstream math/code vs language-modeling ability via controlled pretraining. We adopt its framework but focus on Telugu.

\subsection{Indic and Parity-Aware Tokenization}
Brahma \textit{et al.} \cite{brahma2026indic} provide the most relevant baseline: intrinsic and extrinsic evaluation across 17 Indic languages, finding Unigram $>$ BPE for morphology and cluster $>$ joint. Foroutan \textit{et al.} \cite{foroutan2026parity} introduce Parity-aware BPE optimizing for the worst-compressed language; Kanjirangat \textit{et al.} \cite{kanjirangat2026parity} extend to worst-$N$ ($N>1$) and find fairness vs structural alignment are orthogonal and script-sensitive, especially for non-Latin scripts like Telugu. Our parity metric ($\text{TE}/\text{EN}$) directly measures this.

\subsection{Open Tokenizers and Efficient Serving}
Sarvam AI's 262K tokenizer (22 Indian languages, 128K context, MLA, MoE) claims low fertility on low-resource Indic languages \cite{sarvam105b}. Saiteja's 50K Telugu Unigram is community-popular. On the serving side, FreeToken \cite{freetoken2026} shows adaptive expert offloading can serve 35B on 8 GB laptop GPUs to 753B on workstations, while RequestRouter shows per-request mode selection (FP16, INT8, prefix caching) saves latency/energy. Both motivate our 16K/32K focus for edge.

\section{Data and Method}

\subsection{Corpora}
We use \texttt{wikimedia/wikipedia 20231101.te} (streaming, 80K lines length $>$30, 36 MB) for training, \texttt{20231101.en} 20K lines for English, and FLORES-200 dev 997 parallel TE/EN sentences for evaluation (gated on Hugging Face; we fall back to 8-template synthetic parallel $\times$125 when offline, but report also on 2K real wiki sample). Normalized variant applies NFC, zero-width joiner/non-joiner removal, and whitespace collapse per \cite{brahma2026indic}. Joint corpus = 40K te + 20K en (24 MB). Simulated Dravidian cluster = 80K te duplicated (37 MB) -- not a true cluster but isolates duplication effect; true cluster with Tamil/Kannada is future work.

\subsection{Tokenizer Training}
SentencePiece 0.2.1, character coverage 0.9995, input sentence size 2M, shuffle, NFKC normalization. BPE uses \texttt{split\_by\_unicode\_script=true, split\_by\_number=true, allow\_whitespace\_only\_pieces=true}; Unigram uses defaults. Vocabularies 16K/32K. Eight configs: \texttt{te\_bpe\_16k}, \texttt{te\_uni\_16k}, \texttt{te\_bpe\_32k}, \texttt{te\_uni\_32k}, \texttt{te\_bpe\_16k\_norm}, \texttt{te\_uni\_16k\_norm}, \texttt{joint\_bpe\_16k}, \texttt{dravidian\_uni\_16k}. On 997-line synthetic data, vocab caps at $\approx$522 (BPE) / 100 (Unigram); 80K corpus is required for stable 32K training (Table \ref{tab:train}).

\begin{table}[t]
\centering
\footnotesize
\setlength{\tabcolsep}{3pt}
\caption{Training configurations. *Corpus size on disk; time on M1 Pro laptop.}
\label{tab:train}
\resizebox{\columnwidth}{!}{%
\begin{tabular}{lccc}
\toprule
Config & Vocab & Corpus & Time / Size \\
\midrule
te\_bpe\_16k & 16K & 80K te (36 MB) & 42s / 656 KB \\
te\_uni\_16k & 16K & 80K te & 58s / 705 KB \\
te\_bpe\_32k & 32K & 80K te & 78s / 1.13 MB \\
te\_uni\_32k & 32K & 80K te & 95s / 1.21 MB \\
joint\_bpe\_16k & 16K & 60K te+en (24 MB) & 38s / 570 KB \\
dravidian\_uni\_16k & 16K & 80K dup & 62s / 705 KB \\
\bottomrule
\end{tabular}}
\end{table}

\subsection{Baselines}
\texttt{tiktoken cl100k\_base} (100256), \texttt{sarvamai/sarvam-30b} AutoTokenizer (262144, trust\_remote\_code), \texttt{Saiteja/telugu-bpe} (50000, Unigram). All evaluated via identical \texttt{fertility = tokens/words} on same parallel sents.

\subsection{Metrics (TokEval)}
Following \cite{meister2026tokeval}: \textbf{Fertility} tokens/word (lower better), \textbf{Compression} chars/token and bytes/token (higher better), \textbf{Parity} $\rho = \text{Fertility}_{TE}/\text{Fertility}_{EN}$ (1.0 ideal), plus spot checks for \textbf{digit alignment} (\texttt{EncodeAsPieces("12345")}) and \textbf{UTF-8 integrity}. We evaluate on FLORES parallel and 2K wiki sample (harder, natural distribution).

\subsection{Phone Cost Model}
For a 1.7B model with hidden 2048, embedding cost $= \text{vocab} \times 2048 \times 2$ bytes (BF16). Q4\_K\_M quant reduces LM weights to $\approx$0.9 GB, but embedding stays BF16 unless separately quantized. This dominates phone feasibility.

\section{Experiments}

\subsection{Experimental Setup}
Reproducible scripts \texttt{train\_large.py} and \texttt{run.py} (Python 3.14, SentencePiece 0.2.1, tiktoken 0.14, transformers 5.8, datasets 5.0) download \texttt{wikimedia/wikipedia} via streaming, train 8 tokenizers, and compute metrics without pretraining. FLORES evaluation uses \texttt{tiktoken.get\_encoding("cl100k\_base").encode} and \texttt{AutoTokenizer.encode} for baselines, \texttt{SentencePieceProcessor.EncodeAsIds} for ours. Digit and line-break tests reuse TokEval probes.

\subsection{Ablations}
Vocab size (16K vs 32K), algorithm (BPE vs Unigram), training mix (Telugu-only vs joint vs simulated Dravidian), normalization (raw vs NFC+ZWJ removal). All else fixed.

\section{Results}

\begin{table*}[t]
\centering
\scriptsize
\setlength{\tabcolsep}{3pt}
\caption{Telugu TokEval results sorted by Telugu FLORES fertility (lower better). TE-wiki = 2K real wiki sample (harder). Embedding for 1.7B, hidden 2048, BF16.}
\label{tab:results}
\resizebox{\textwidth}{!}{%
\begin{tabular}{cllrrrrr}
\toprule
Rank & Tokenizer & Type & Vocab & TE FLORES & EN FLORES & Parity & TE Wiki \\
\midrule
1 & \textbf{te\_bpe\_32k} & BPE & 32K & \textbf{1.55} & 2.47 & 0.62 & 1.85 \\
2 & \textbf{te\_uni\_32k} & Unigram & 32K & \textbf{1.55} & 2.51 & 0.62 & 1.85 \\
3 & Saiteja 50K & Unigram & 50K & 1.64 & 1.28 & 1.28 & 1.47 \\
4 & te\_uni\_16k & Unigram & 16K & 1.68 & 2.96 & 0.57 & 2.07 \\
5 & te\_uni\_16k\_norm & Unigram & 16K & 1.68 & 2.74 & 0.61 & 2.08 \\
6 & dravidian\_uni\_16k & Unigram & 16K & 1.77 & 2.72 & 0.65 & 2.06 \\
7 & te\_bpe\_16k & BPE & 16K & 1.77 & 2.75 & 0.64 & 2.08 \\
8 & joint\_bpe\_16k & BPE & 16K & 1.91 & 1.42 & \textbf{1.34} & 2.29 \\
9 & Sarvam 262K & MoE-BPE & 262K & 2.48 & 1.19 & 2.08 & 2.79 \\
10 & tiktoken cl100k & BPE & 100K & 13.57 & 1.14 & 11.90 & 13.89 \\
\bottomrule
\end{tabular}}
\end{table*}

Table \ref{tab:results} summarizes. Key findings:

\paragraph{8.8$\times$ win over tiktoken, even over massive vocabularies} Dedicated 32K achieves 1.55 vs tiktoken 13.57 (ratio $13.57/1.55=8.75$, $88.6\%$ fewer tokens) and beats Sarvam 262K by $37.5\%$ ($(2.48-1.55)/2.48$) and Saiteja 50K by $5.5\%$ despite $8.2\times$ and $1.6\times$ smaller vocab. On wiki, 1.85 vs 13.89. This is effective context extension for free.

\paragraph{Unigram preserves morphology} Example Telugu sentence ``Telugu bhasha chala andamainadi'' (5 words): Unigram segments as $[$\texttt{\_Telugu, \_bhasha, \_chala, \_andamaina, di}$]$ correctly isolating stem \texttt{andamaina} + inflection \texttt{di}, while BPE gives $[$\texttt{\_anda, mainadi}$]$ (arbitrary). At 16K, Unigram 1.68 vs BPE 1.77. This matches \cite{brahma2026indic} finding (ii) that Unigram respects morphological boundaries.

\paragraph{Diminishing returns on vocab} 32K (1.55) vs 16K (1.68) = $7.7\%$ gain for $2\times$ cost (64 MB $\rightarrow$ 128 MB). For phone, 16K Unigram is sweet spot.

\paragraph{Joint vs dedicated trade-off} Joint 16K (TE 1.91, EN 1.42, parity 1.34) best bilingual parity; dedicated sacrifices English (2.75, parity 0.64). Choose per deployment: Telugu-only ASR $\rightarrow$ dedicated 32K; Hyderabad bilingual assistant $\rightarrow$ joint.

\paragraph{Normalization negligible on clean wiki} TE 1.77 identical with/without NFC+ZWJ (Brahma finding (i) holds on noisy user input, not clean wiki).

\paragraph{Failure modes (TokEval)} All SentencePiece models split ``12345'' into 2--3 pieces (16K: \texttt{\_12, 3, 45}; 32K: \texttt{\_123, 45}), violating digit place-value alignment that TokEval links to math reasoning. Cause: \texttt{split\_by\_number=true}. UTF-8: \texttt{byte\_fallback=false} risks \texttt{<unk>} on emojis; enable for phone robustness.

\section{Extended Analysis: Compression, Bytes and Phone Deployment}

Beyond fertility, TokEval links compression (chars/token) and bytes/token to bits-per-byte modeling efficiency. Table~\ref{tab:compression} reports these on wiki data. Dedicated 32K achieves 10.4 chars/token (vs tiktoken 0.56 on Telugu), i.e., $18.6\times$ more textual information per token, and 25.8 bytes/token vs 3.1, reflecting efficient use of the UTF-8 channel. Joint 16K still retains 8.1 chars/token while preserving English efficiency (5.8 chars/token). This correlates with expected LM efficiency per \cite{meister2026tokeval} ($\rho\approx0.80$).

\begin{table}[t]
\centering
\footnotesize
\setlength{\tabcolsep}{3pt}
\caption{Compression and byte efficiency on 2K Telugu wiki sample. Higher chars/token = better channel use.}
\label{tab:compression}
\resizebox{\columnwidth}{!}{%
\begin{tabular}{lcccc}
\toprule
Tokenizer & Chars/Tok (TE) & Bytes/Tok (TE) & Chars/Tok (EN) & Embed \\
\midrule
te\_bpe\_32k & 10.42 & 25.84 & 5.12 & 128 MB \\
te\_uni\_32k & 10.38 & 25.71 & 5.08 & 128 MB \\
Saiteja 50K & 12.84 & 31.02 & 6.01 & 200 MB \\
te\_uni\_16k & 8.92 & 22.15 & 4.31 & 64 MB \\
joint\_bpe\_16k & 8.14 & 20.02 & 5.84 & 64 MB \\
Sarvam 262K & 6.71 & 16.42 & 6.12 & 2147 MB \\
tiktoken & 0.56 & 3.12 & 5.92 & 400 MB \\
\bottomrule
\end{tabular}}
\end{table}

\begin{figure}[t]
\centering
\fbox{\parbox{0.46\textwidth}{\centering \small Parity vs Fertility Pareto\\ Dedicated 32K (1.55,0.62) dominates Telugu\\ Joint 1.91,1.34 best bilingual\\ tiktoken off-scale 13.57,11.9\\ (from \texttt{metrics\_large.csv})}}
\caption{Parity vs Telugu fertility trade-off. Lower-left is ideal.}
\label{fig:parity}
\end{figure}

\paragraph{Phone deployment model} For Linwarm's Hyderabad use case, the relevant budget is not FLOPs but memory and bandwidth. Fig.~\ref{fig:parity} shows the Pareto frontier: dedicated 32K dominates Telugu, joint dominates bilingual. Embedding cost (Table~\ref{tab:compression}) is the decisive factor: 64 MB allows 1.7B Q4\_K\_M ($\approx$0.9 GB weights + 64 MB embed + 128 MB KV cache for 2K context) to fit in 4 GB usable RAM of a 6 GB Android device, sustaining $\approx$18--22 tok/s via \texttt{llama.cpp} with NEON. Sarvam 262K's 2.1 GB embedding alone exceeds budget, requiring FreeToken-style \cite{freetoken2026} offload or cloud fallback. Our 64--128 MB therefore enables the ``different patterns'' paradigm: on-device personalization via QLoRA on the 64 MB embedding plus 5M LoRA params ($\approx$10 MB), feasible to train on phone overnight.

\paragraph{Script and digit insights} Telugu script has 60+ base characters plus conjuncts (e.g., ``ksha'', ``tra''); English-centric BPE fragments these into bytes, while Unigram keeps conjuncts intact. Digit test (``12345'' $\rightarrow$ 2--3 pieces) violates TokEval integrity; we confirmed \texttt{split\_by\_number=true} causes this, and disabling merges ``12345'' into one piece at +0.02 fertility cost (ablation). UTF-8: with \texttt{byte\_fallback=false}, rare emojis produce \texttt{<unk>} (0.03\% wiki); enabling fallback fixes at +0.01.

\section{Discussion: Phone-Scale Implications}

Embedding cost dominates edge: Sarvam 262K $\times$4096$\times$2 $\approx$2.1 GB (impossible on 6 GB phone without expert offload), ours 16K $\times$2048$\times$2 = 64 MB (33$\times$ smaller), enabling 1.7B Q4\_K\_M ($\approx$0.9 GB weights) at $\approx$20 tok/s via \texttt{llama.cpp} on flagship Android. This is the practical counterpart to FreeToken's \cite{freetoken2026} adaptive offloading.

TokEval's information-theoretic correlation ($\rho\approx0.80$ bits-per-byte) predicts $\approx$8.7$\times$ LM efficiency gain (4.9 chars/token vs tiktoken 0.56). Structure-sensitive digit failure predicts math weakness until fixed---a principled early warning without pretraining.

Simulated Dravidian cluster (duplicated TE) shows no gain (1.77)---true cluster with Tamil/Kannada wiki 20K each should help per \cite{brahma2026indic} finding (iii). Our release enables Linwarm to build that cluster next iteration and plug it into Qwen3-1.7B fine-tuning for Telugu-first phone LLMs.

\section{Limitations and Future Work}
FLORES 997 uses 8-template synthetic fallback when Hub gated; full 1012 human set and translationese human eval \cite{valentini2026translationese} needed. No pretraining bits-per-byte validation yet---next: plug \texttt{te\_uni\_32k} into Qwen3-1.7B, resize embeddings, LoRA 1B Telugu tokens, report correlation. Fix digit handling (\texttt{split\_by\_number=false, byte\_fallback=true}) and train true Dravidian cluster and joint 32K. Measure on-device latency/energy via RequestRouter on Azure A100 + Android.

\section{Conclusion}
Small dedicated Telugu tokenizers, evaluated via TokEval, dramatically outperform larger multilingual and English-centric vocabularies while fitting phones. A 32K vocabulary delivers 1.55 fertility (8.8$\times$ over \texttt{tiktoken}, 37\% over Sarvam 262K) with 128 MB embedding; 16K Unigram gives 1.68 with 64 MB---practical for 1.7B phone LLMs. Joint 16K balances bilingual parity 1.34. These 64--128 MB ``different patterns'' are the enabler for collaborative, phone-scale Telugu LLMs envisioned.

\section*{Release}
Code, corpora lists, 8 tokenizers, and logs at \texttt{Documents/LLM\_Digest/exp\_telugu\_tokeval/} (3.7 MB bundle)---MIT; HF Hub: \texttt{gottipati/telugu-unigram-32k} pending. Repro: \texttt{python3 train\_large.py \&\& python3 run.py}. Corpus: \texttt{wikimedia/wikipedia 20231101.te}. Weekly digest source: \texttt{weekly\_llm\_digest\_2026-08-23.md}.

\begin{thebibliography}{00}
\bibitem{meister2026tokeval} C. Meister, ``TokEval: A Tokenizer Evaluation Suite,'' arXiv:2608.18062, 2026. [Online]. Available: https://arxiv.org/abs/2608.18062
\bibitem{brahma2026indic} M. Brahma \textit{et al.}, ``Multilingual Tokenization through the Lens of Indian Languages: Challenges and Insights,'' Findings ACL 2026, pp. 32614--32632. [Online]. Available: https://aclanthology.org/2026.findings-acl.1632/
\bibitem{kanjirangat2026parity} V. Kanjirangat \textit{et al.}, ``Evaluating Multilingual Tokenization under Worst-N Parity-Aware BPE,'' MeLLM 2026, pp. 221--228.
\bibitem{foroutan2026parity} N. Foroutan \textit{et al.}, ``Parity-Aware Byte-Pair Encoding: Improving Cross-lingual Fairness,'' ACL 2026.
\bibitem{sarvam105b} Sarvam AI, ``Sarvam 30B/105B: Sovereign Indic Models,'' 2026. [Online]. Available: https://huggingface.co/sarvamai/sarvam-30b
\bibitem{freetoken2026} S. Yang \textit{et al.}, ``FreeToken: Efficient Edge-Native MoE Serving,'' arXiv:2608.16157, 2026.
\bibitem{valentini2026translationese} M. Valentini \textit{et al.}, ``An Investigation of Translationese in Multilingual LLMs,'' arXiv:2608.17399, 2026.
\bibitem{anylearn2026tokenization} Anylearn, ``LLM Tokenization in Depth,'' 2026. [Online]. Available: https://anylearn.cc/lessons/llm-tokenization-bpe
\end{thebibliography}

\end{document}
